
Hilbert spaces are a special type of a more general class of spaces
known as Banach spaces. We are interested in Banach spaces not just
for the sake generality, but also because they naturally appear in
Hilbert space theory. For instance, the space of bounded linear maps
on a Hilbert space is not itself a Hilbert space, but only a Banach
space.

\section{Generalities on Banach spaces}

We begin with some basis notions from metric space theory.

\begin{definition}
A \emph{metric space}\index{metric space} is a pair $(X,d)$, where $X$
is a set and $d$ is a \emph{metric} on $X$, that is, a map $d\cl
X\times X \to [0,\infty)$ satisfying
\begin{enumerate}[label=(\roman*)]
\item $d(x,y) = 0\ \Leftrightarrow \ x=y$ \hfill (identity of indiscernibles)
\item $d(x,y)=d(y,x)$ \hfill (symmetry)
\item $d(x,z)\leq d(x,y)+d(y,z)$ \hfill (triangle inequality)
\end{enumerate}
for all $x,y,z\in X$.
\end{definition}

\begin{definition}
A sequence $\{x_n\}_{n\in \N}$ in a metric space $(X,d)$ is said to
\emph{converge}\index{convergence} to an element $x\in X$, written
$\displaystyle \lim_{n\to \infty}x_n=x$, if
\begin{equation*}
\forall \, \varepsilon > 0 : \exists \, N \in \N : \forall \, n \geq N : \ d(x_n,x)<\varepsilon.
\end{equation*}
\end{definition}

A sequence in a metric space can converge to at most one element.

\begin{definition}
A \emph{Cauchy sequence}\index{Cauchy sequence} in a metric space $(X,d)$ is a sequence $\{x_n\}_{n\in \N}$ such that
\begin{equation*}
\forall \, \varepsilon > 0 : \exists \, N \in \N : \forall \, n,m \geq N : \ d(x_n,x_m)<\varepsilon.
\end{equation*}
\end{definition}
Any convergent sequence is clearly a Cauchy sequence.

\begin{definition}
A metric space $(X,d)$ is said to be \emph{complete}\index{completeness} if every Cauchy sequence converges to some $x\in X$.
\end{definition}
A natural metric on a vector space is that induced by a norm.
\begin{definition}
A \emph{normed space}\index{normed space} is a (complex) vector space
$(V,+,\cdot)$ equipped with a \emph{norm}, that is, a map
$\|\cdot\|\cl V \to [0,\infty)$ satisfying
\begin{enumerate}[label=(\roman*)]
\item $\|f\| = 0 \ \Leftrightarrow\ f=0$ \hfill (positive definiteness)
\item $\|z\cdot f\|=|z|\|f\|$ \hfill (homogeneity/scalability)
\item $\|f+g\|\leq \|f\|+\|g\|$ \hfill (triangle inequality/sub-additivity)
\end{enumerate}
for all $f,g\in V$ and all $z\in \C$.
\end{definition}
One we have a norm $\|\cdot\|$ on $V$, we can define a metric $d$ on $V$ by 
\begin{equation*}
d(f,g) := \|f-g\|.
\end{equation*}
Then we say that the normed space $(V,\|\cdot\|)$ is \emph{complete} if the metric space $(V,d)$, where $d$ is the metric induced by $\|\cdot\|$, is complete. Note that we will usually suppress inessential information in the notation, for example writing $(V,\|\cdot\|)$ instead of $(V,+,\cdot,\|\cdot\|)$.

\begin{definition}
A \emph{Banach space}\index{Banach space} is a complete normed vector space.
\end{definition}

\begin{example}
  The set of continuous functions
  \begin{equation}
    C^0_\C[0,1]:=\{f\cl [0,1]\to \C \mid f \text{ is continuous}\}
  \end{equation}
  is a vector space with the pointwise sum and scalar product (meaning
  that the sum and product of continuous functions is again continuous).
  Since $[0,1]$ is closed and bounded, it is compact and hence
  every complex-valued continuous function $f\cl [0,1]\to \C$ is
  bounded, in the sense that the supremum
  \begin{equation}
    \|f\|_\infty = \sup \{|f(x)|\mid x\in[0,1]\}
  \end{equation}
  is always finite.
  We claim that this function defines a norm on $C^0_\C[0,1]$,
  called the \emph{supremum} (or \emph{infinity}) \emph{norm},
  and that, equipped with this norm, the space \(C^{0}_\C[0,1]\) is
  complete (i.e. it is a Banach space).
  Let us show this in some detail.

  The function \(\|.\|_\infty\) is clearly positive definite, so to
  see that it is a norm, it remains to show homogeneity and the
  triangle inequality.
  Let $f,g\in C^0_\C[0,1]$ and $z\in \C$. Then we have
  \begin{align}
    \|z\cdot f\|_{\infty}
    &:= \sup_{x\in[0,1]}|zf(x)| \\
    &= \sup_{x\in[0,1]}|z||f(x)| \\
    & =|z|\sup_{x\in[0,1]}|f(x)| \\
    &= |z|\|f\|_{\infty}
  \end{align}
  and
  \begin{align*}
  \|f+g\|_{\infty}
  &:= \sup_{x\in[0,1]}|(f+g)(x)|\\
  &= \sup_{x\in[0,1]}|f(x)+g(x)|\\
  &\leq \sup_{x\in[0,1]}(|f(x)|+|g(x)|)\\
  & = \sup_{x\in[0,1]}|f(x)|\,+\sup_{x\in[0,1]}|g(x)|\\
  & = \|f\|_{\infty}+\|g\|_{\infty}.
  \end{align*}
  Hence, $(C^0_\C[0,1],\|\cdot\|_{\infty})$ is indeed a normed space.
  We now show that $C^0_\C[0,1]$ is complete. Let $\{f_n\}_{n\in \N}$ be
  a Cauchy sequence of functions in $C^0_\C[0,1]$, that is
  \begin{equation*}
    \forall \, \varepsilon > 0 :
    \exists \, N \in \N :
    \forall \, n,m \geq N :
    \ \|f_n-f_m\|_{\infty} <\varepsilon.
  \end{equation*}
  We seek an $f\in C^0_\C[0,1]$ such that $\displaystyle
  \lim_{n\to\infty}f_n=f$. We will proceed in three steps.
  \begin{enumerate}
    \item
      We first define a candidate function \(f\).
      Fix $y\in [0,1]$ and $\varepsilon >0$. By definition of supremum, we have
      \begin{equation*}
      |f_n(y) -f_m(y)|
      \leq \sup_{x\in[0,1]}|f_n(x)-f_m(x)|
      =: \|f_n-f_m\|_{\infty}.
      \end{equation*}
      Hence, there exists $N\in \N$ such that
      \begin{equation*}
      \forall \, n,m \geq N : \ |f_n(y) -f_m(y)|<\varepsilon,
      \end{equation*}
      that is, the sequence of complex numbers $\{f_n(y)\}_{n\in \N}$ is a
      Cauchy sequence. Since $\C$ is a complete metric space,
      the sequence \(f_n(y)\) converges. Thus, we can define a function
      \begin{align*}
      f\cl [0,1] & \to \C
      \\
      x & \mapsto \lim_{n\to\infty}f_n(x),
      \end{align*}
      called the \emph{pointwise limit} of $f$.
      Note that, even if \(f(x)=\lim_{n\to\infty}f_n(x)\), this does
      \emph{not} automatically imply that $\lim_{n\to
      \infty}f_n=f$ (convergence with respect to the supremum norm), nor that
      $f\in C^0_\C[0,1]$, and hence we need to check separately that these
      do, in fact, hold.
    \item
      Let us check that $f$ is continuous, i.e.
      that for every \(\epsilon>0\), there is some \(\delta>0\) such that
      \begin{equation}
        |x-y|<\delta \implies |f(x)-f(y)|<\epsilon
      .\end{equation}
      For any \(n\in \N\), we have
      \begin{align*}
      |f(x)-f(y)|
      & = |f(x)-f_n(x)+f_n(x)-f_n(y)+f_n(y)-f(y)|
      \\
      & \leq |f(x)-f_n(x)|+|f_n(x)-f_n(y)|+|f_n(y)-f(y)|.
      \end{align*}
      Note that the summand in the middle can be bound by \(\epsilon /
      3\) by taking \(\delta\) small enough.
      The main challenge is that the other two summands are bounded
      separately for each point, e.g. for each \(x\) there is \(N_x\)
      such that
      \begin{equation}
        |f(x)-f_{N_x}(x)|<\frac{\epsilon}{3}
      .\end{equation}
      If we could find a single \(N\) that works for all \(x\) at
      once, we would have
      \begin{align*}
      |f(x)-f(y)|
      & \leq \frac{\epsilon}{3} + |f_N(x)-f_N(y)|+\frac{\epsilon}{3}
      \end{align*}
      so by continuity of \(f_N\) there is a \(\delta>0\) such that
      \(|f_N(x)-f_N(y)|<\epsilon / 3\), wich implies the result.

      It is the Cauchy property of \(\{f_n\}\) that comes to save the day.
      There is an \(N\in\N\) such that
      \begin{equation}
        \forall x\in[0,1], \ \forall m,n\geq N,\ |f_m(x)-f_n(x)| <
        \epsilon / 4
      \end{equation}
      Taking \(n=N\) and \(m\to\infty\), we obtain
      \begin{equation}
        \forall x\in[0,1], \ |f(x)-f_N(x)|\leq \epsilon / 4 < \epsilon /
        3,
      \end{equation}
      which is what we needed.
    \item
      Finally, it remains to show that $\lim_{n\to \infty}f_n=f$ or,
      equivalently, \(\lim_{n\to\infty}\|f-f_n\|=0\).
      Given \(\epsilon>0\), we must show that there is \(N\) such that
      \begin{equation}
        \forall n\geq N,\ \|f-f_n\| < \epsilon
      .\end{equation}
      The Cauchy property of \(\{f_n\}\) tells us
      that there is \(N\geq 0\) with
      \begin{equation}
        \forall m,n\geq N,\ \|f_m-f_n\|<\epsilon / 2
      .\end{equation}
      This implies that
      \begin{equation}
        \forall x\in[0,1],\ \forall m,n\geq N,\ |f_m(x)-f_n(x)|<\epsilon
        / 2
      .\end{equation}
      Taking the limit \(m\to\infty\), we get
      \begin{equation}
        \forall x\in[0,1],\ \forall n\geq N,\ |f(x)-f_n(x)|\leq \epsilon
        / 2
      \end{equation}
      so, for all \(n\geq N\), we have \(\|f-f_n\|\leq
      \frac{\epsilon}{2}<\epsilon \), as we wanted.
  \end{enumerate}
  This completes the proof that $(C^0_\C[0,1],\|\cdot\|_{\infty})$ is a
  Banach space.
\end{example}

\begin{remark}
The previous example shows that checking that something is a Banach space, and the completeness property in particular, can be quite tedious. However, in the following, we will typically already be working with a Banach (or Hilbert) space and hence, rather than having to check that the completeness property holds, we will instead be able to use it to infer the existence (within that space) of the limit of any Cauchy sequence.
\end{remark}

\section{Bounded linear operators}

As usual in mathematics, once we introduce a new types of structure,
we also want study maps between instances of those structures, with
extra emphasis placed on the structure-preserving maps.
We begin with linear maps between normed spaces.

\begin{proposition}
\label{prp:boundequiv}
Let $(V,\|\cdot\|_V)$, $(W,\|\cdot\|_W)$ be normed spaces
and $A\cl V\to W$ a linear operator.
The following conditions are equivalent.
\begin{enumerate}[label=(\roman*)]
  \item \(\displaystyle \sup_{f\in V-\{0\}}\frac{\|Af\|_W }{ \|f\|_V}
    < \infty\)
  \item $\displaystyle \sup_{\|f\|_V=1}\|Af\|_W< \infty$
  \item $\exists \, k > 0 : \forall \, f\in V: \ \|f\|_V \leq 1\, \Rightarrow \, \|Af\|_W \leq k$
  \item $\exists \, k > 0 : \forall \, f\in V: \ \|Af\|_W\leq  k \|f\|_V$
  \item the map $A\cl V \to W$ is continuous with respect to the topologies induced by the respective norms on $V$ and $W$
  \item the map $A$ is continuous at $0\in V$. 
\end{enumerate}
Moreover, in this case, the following numbers are equal:
\begin{itemize}
  \item \(\sup\left\{\frac{\|Af\|}{\|f\|} \mid f\in V-\{0\}\right\}\)
  \item \(\sup \{\|Af\| : f\in V, \|f\|=1\}\)
  \item \(\sup \{\|Af\| : f\in V, \|f\|\leq 1\}\)
  \item \(\inf \{C\in[0,\infty) \mid \forall f\in V, \|Af\|\leq C\|f\|\}\)
\end{itemize}
\end{proposition}

\begin{definition}
  Let \(V,W\) be normed spaces.
  A linear operator $A\cl V\to W$ is said to be \emph{bounded} if it
  satisfies the conditions in the previous proposition.
  If $A\cl V \to W$ is a bounded operator, the \emph{operator
  norm}\index{operator norm} \(\|A\|\) of $A$ is the number given by
  any of the three suprema or the infimum in the proposition.
\end{definition}

\begin{remark}
  The set of \(C\in[0,\infty)\) with \(\|Af\|\leq C\|f\|\) for all
  \(f\in V\) is a closed set in \(\R\).
  Indeed, if \(C_n\) is a sequence in the set, we have \(\|Af\|\leq
  C_n\|f\|\), so \(\|Af\|\leq\lim_{n\to\infty}C_n\|f\|\).
  Then, the operator norm \(\|A\|\) is actually the \emph{minimum} of
  this set, meaning that
  \begin{equation}
    \forall f\in V, \|Af\|\leq\|A\|\|f\|
  .\end{equation}
\end{remark}

\begin{example}
Let $\id_W\cl W\to W$ be the identity operator on a Banach space $W$. Then
\begin{equation*}
\sup_{f\in W}\frac{\|\id_Wf\|_W}{\|f\|_W} =\sup_{f\in W}1 = 1 <\infty.
\end{equation*}
Hence, $\id_W$ is a bounded operator and has unit norm.
\end{example}

\begin{example}
Denote by $C^1_{\C}[0,1]$ the complex vector space of once continuously differentiable complex-valued functions on $[0,1]$. Since differentiability implies continuity, this is a subset of $C^0_{\C}[0,1]$. Moreover, since sums and scaling by a complex number of continuously differentiable functions are again continuously differentiable, this is, in fact, a vector subspace of $C^0_{\C}[0,1]$, and hence also a normed space with the supremum norm $\|\cdot\|_{\infty}$.

Consider the first derivative operator
\begin{align*}
D \cl C^1_{\C}[0,1] & \to C^0_{\C}[0,1]\\
 f & \mapsto f'.
\end{align*}
We know from undergraduate real analysis that $D$ is a linear operator. We will now show that $D$ is an unbounded\footnote{Some people take the term unbounded to mean ``not necessarily bounded''. We take it to mean ``definitely not bounded'' instead.} linear operator. That is,
\begin{equation*}
\sup_{f\in C^1_{\C}[0,1]} \frac{\|Df\|_{\infty}}{\|f\|_{\infty}} = \infty.
\end{equation*}
Note that, since the norm is a function into the real numbers, both $\|Df\|_{\infty}$ and $\|f\|_{\infty}$ are always finite for any $f\in C^1_{\C}[0,1]$. Recall that the supremum of a set of real numbers is its least upper bound and, in particular, it need not be an element of the set itself. What we have to show is that the set
\begin{equation*}
\biggl\{ \frac{\|Df\|_{\infty}}{\|f\|_{\infty}} \ \Big| \ f\in C^1_{\C}[0,1] \biggr\} \subset \R
\end{equation*}
contains arbitrarily large elements. One way to do this is to exhibit a positively divergent (or unbounded from above) sequence within the set. 

Consider the sequence $\{f_n\}_{n\geq 1}$ where $f_n(x):=\sin(2\pi nx)$. We know that sine is continuously differentiable, hence $f_n\in C^1_{\C}[0,1]$ for each $n\geq 1$, with
\begin{equation*}
Df_n(x) = D(\sin(2\pi nx)) = 2\pi n\cos(2\pi nx).
\end{equation*}
We have
\begin{equation*}
\|f_n\|_{\infty} = \sup_{x\in[0,1]}|f_n(x)| = \sup_{x\in[0,1]}|\sin(2\pi nx)| = \sup \, [-1,1] = 1 
\end{equation*}
and
\begin{equation*}
\|Df_n\|_{\infty} = \sup_{x\in[0,1]}|Df_n(x)| = \sup_{x\in[0,1]}|2\pi n\cos(2\pi nx)| = \sup\, [-2\pi n,2\pi n] = 2\pi n.
\end{equation*}
Hence, we have
\begin{equation*}
\sup_{f\in C^1_{\C}[0,1]} \frac{\|Df\|_{\infty}}{\|f\|_{\infty}} \geq \sup_{\{f_n\}_{n\geq 1}} \frac{\|Df\|_{\infty}}{\|f\|_{\infty}} = \sup_{n\geq 1}\,2\pi n = \infty,
\end{equation*}
which is what we wanted. As an aside, we note that $C^1_{\C}[0,1]$ is not complete with respect to the supremum norm, but it is complete with respect to the norm
\begin{equation*}
\|f\|_{C^1}:=\|f\|_{\infty}+\|f'\|_{\infty}.
\end{equation*}
While the derivative operator is still unbounded with respect to this new norm, in general, the boundedness of a linear operator does depend on the choice of norms on its domain and target, as does the numerical value of the operator norm.
\end{example}

\begin{remark}
Apart from the ``minor'' detail that in quantum mechanics we deal with Hilbert spaces, use a different norm than the supremum norm and that the (one-dimensional) momentum operator acts as $P(\psi):=-\mathrm{i}\hbar\psi'$, the previous example is a harbinger of the fact that the momentum operator in quantum mechanics is unbounded. This will be the case for the position operator $Q$ as well.
\end{remark}

\begin{lemma}
\label{lem:forcompl}
Let $(V,\|\cdot\|)$ be a normed space. Then, addition, scalar
multiplication, and the norm are all sequentially continuous. That is,
for any sequences $\{f_n\}_{n\in \N}$ and $\{g_n\}_{n\in \N}$ in $V$
converging to $f\in V$ and $g\in V$ respectively, and any sequence
$\{z_n\}_{n\in \N}$ in $\C$ converging to $z\in \C$, we have
\begin{enumerate}[label=(\roman*)]
\item $\displaystyle \lim_{n\to \infty}(f_n+g_n)=f+g$
\item $\displaystyle \lim_{n\to \infty}z_nf_n=zf$.
\item $\displaystyle \lim_{n\to \infty}\|f_n\|=\|f\|$
\end{enumerate}
\end{lemma}

\iffalse
\begin{proof}
\begin{enumerate}[label=(\roman*)]
\item Let $\varepsilon >0$. Since $\displaystyle \lim_{n\to \infty}f_n=f$ and $\displaystyle \lim_{n\to \infty}g_n=g$ by assumption, there exist $N_1,N_2\in \N$ such that
\begin{align*}
  \forall \, n\geq N_1, & \ \|f-f_n\|<\tfrac{\varepsilon}{2}\phantom{.}\\
  \forall \, n\geq N_2, & \ \|g-g_n\|<\tfrac{\varepsilon}{2}.
\end{align*}
Let $N:=\max\{N_1,N_2\}$. Then, for all $n\geq N$, we have
\begin{align*}
\|(f_n+g_n)-(f+g)\|
&= \|f_n-f+g_n-g\|\\
&\leq \|f_n-f\|+\|g_n-g\|\\
& < \tfrac{\varepsilon}{2}+\tfrac{\varepsilon}{2}\\
& = \varepsilon.
\end{align*}
Hence $\displaystyle \lim_{n\to \infty}(f_n+g_n)=f+g$.


\item Since $\{z_n\}_{n\in \N}$ is a convergent sequence in $\C$, it is bounded. That is,
\begin{equation*}
\exists \, k>0 : \forall \, n\in \N : \ |z_n|\leq k.
\end{equation*}
Let $\varepsilon >0$. Since $\displaystyle \lim_{n\to \infty}f_n=f$ and $\displaystyle \lim_{n\to \infty}z_n=z$, there exist $N_1,N_2\in \N$ such that
\begin{align*}
  \forall \, n\geq N_1, & \ \|f-f_n\|<\frac{\varepsilon}{2k}\\
  \forall \, n\geq N_2, & \ \|z-z_n\|<\frac{\varepsilon}{2\|f\|}.
\end{align*}
Let $N:=\max\{N_1,N_2\}$. Then, for all $n\geq N$, we have
\begin{align*}
\|z_nf_n-zf\|
&= \|z_nf_n-z_nf+z_nf-zf\|\\
&= \|z_n(f_n-f)+(z_n-z)f\|\\
&\leq \|z_n(f_n-f)\|+\|(z_n-z)f\|\\
&= |z_n|\|f_n-f\|+|z_n-z|\|f\|\\
& < k\frac{\varepsilon}{2k}+\frac{\varepsilon}{2\|f\|}\|f\|\\
& = \varepsilon.
\end{align*}
Hence $\displaystyle \lim_{n\to \infty}z_nf_n=zf$.
\item Let $\varepsilon >0$. Since  $\displaystyle \lim_{n\to \infty}f_n=f$, there exists $N\in \N$ such that
\begin{equation*}
\forall \, n \geq N : \ \|f_n-f\|<\varepsilon.
\end{equation*}
By the triangle inequality, we have
\begin{equation*}
\|f_n\| = \|f_n-f+f\| \leq \|f_n-f\|+\|f\|  
\end{equation*}
so that $\|f_n\|-\|f\|\leq \|f_n-f\|$. Similarly, $\|f\|-\|f_n\|\leq \|f-f_n\|$. Since
\begin{equation*}
\|f-f_n\| = \|-(f_n-f)\| = |-1|\|f_n-f\| = \|f_n-f\|,
\end{equation*}
we have $-\|f_n-f\|\leq\|f_n\|-\|f\|\leq\|f_n-f\|$ or, by using the modulus,
\begin{equation*}
\bigl| \|f_n\|-\|f\|\bigr| \leq \|f_n-f\|.
\end{equation*}
Hence, for all $n\geq N$, we have $\bigl| \|f_n\|-\|f\|\bigr| <\varepsilon$ and thus  $\displaystyle \lim_{n\to \infty}\|f_n\|=\|f\|$.
\qedhere
\end{enumerate}
\end{proof}
\fi
Note that by taking $\{z_n\}_{n\in\N}$ to be the constant sequence whose terms are all equal to some fixed $z\in \C$, we have $\displaystyle \lim_{n\to \infty}zf_n=zf$ as a special case of (ii).

This lemma will take care of some of the technicalities involved in proving the following crucially important result.

\begin{theorem}
Let \(V,W\) be normed spaces. Then the set $\mathcal{L}(V,W)$ of
bounded linear operators, equipped with the operator norm, is a normed space.
Moreover, if \(W\) is Banach, then \(\Cal L(V,W)\) is Banach.
\end{theorem}

\begin{proof}
Define addition and scalar multiplication on $\mathcal{L}(V,W)$ by
\begin{equation*}
(A+B)f := Af+Bf \quad \qquad (zA)f := zAf.
\end{equation*}
It is clear that both $A+B$ and $zA$ are linear operators, so this
addition and scalar multiplication are operations
\begin{align*}
  +\cl \mathcal{L}(V,W)\times \mathcal{L}(V,W) & \to\mathcal{L}(V,W)
  &
  \cdot \cl \C \times \mathcal{L}(V,W) & \to \mathcal{L}(V,W)\\
  (A,B) & \mapsto A+B 
        &
  (z,A) & \mapsto zA
\end{align*}
and one checks that they give \(\Cal L(V,W)\) the structure
of a vector space.
Now we will see that it is a Banach space with the operator norm.
We start by checking that the ``operator norm'' \emph{is} a norm on \(\Cal
L(V,W)\).
\begin{enumerate}[label=(\alph*)]
\item
  We have
  \begin{align*}
    \|(A+B)f\|_W
    &:= \|Af+Bf\|_W \\
    &\leq \|Af\|_W+\|Bf\|_W \\
    &\leq \|A\|\|f\| + \|B\|\|f\| \\
    &= (\|A\|+ \|B\|)\|f\|
  \end{align*}
  since $A$ and $B$ are bounded. Hence, $A+B$ is also bounded and we have
  \begin{equation*}
  \|A+B\|\leq \|A\|+\|B\|.
  \end{equation*}
  \item
  Similarly, for $z\in\C$ and \(A\in\Cal L(V,W)\), we have 
  \begin{align*}
    \|(zA)f\|_W
    &:= \|zAf\|_W \\
    &= |z|\|Af\|_W \\
    &\leq |z|\|A\|\|f\|_W
  \end{align*}
  since $A$ is bounded and $|z|$ is finite.
  Hence, $zA$ is bounded.  and we have
  \begin{equation*}
  \|zA\|=|z|\|A\|.
  \end{equation*}
  \item
    $\displaystyle \|A\|:=\sup_{f\in
    V}\frac{\|Af\|_W}{\|f\|_V}\geq 0$ since $\|\cdot\|_V$ and
    $\|\cdot\|_W$ are norms.
  \item Again, by using the fat that $\|\cdot\|_W$ is a norm,
  \begin{align*}
  \|A\|=0
  & \iff \sup_{f\in V}\frac{\|Af\|_W}{\|f\|_V}= 0 \\
  & \iff \forall \, f\in V:\ \|Af\|_W=0\\
  & \iff \forall \, f\in V:\ Af=0\\
  & \iff A=0.
  \end{align*}
Hence, $(\mathcal{L}(V,W),\|\cdot\|)$ is a normed space.
\end{enumerate}
Now suppose that \(W\) is complete.
We will show that $\Cal{L}(V,W)$ is also complete.
We will proceed in three steps, analogously to the case of $C_{\C}^0[0,1]$.
\begin{enumerate}[label=(c.\roman*)]
\item Let $\{A_n\}_{n\in \N}$ be a Cauchy sequence in
  $\mathcal{L}(V,W)$. Fix $f\in V$ and let $\varepsilon >0$. Then,
  there exists $N\in \N$ such that
  \begin{equation*}
  \forall \, n,m\geq N : \ \|A_n-A_m\|<\frac{\varepsilon}{\|f\|_V}.
  \end{equation*}
  Then, for all $n,m\geq N$, we have
  \begin{equation*}
  \|A_nf-A_mf\|_W
  \leq \|f\|_V\|A_n-A_m\|\\
  < \varepsilon.
  \end{equation*}
  (Note that if $f=0$, we simply have $\|A_nf-A_mf\|_W=0<\varepsilon$
  and, in the future, we will not mention this case explicitly.) Hence,
  the sequence $\{A_nf\}_{n\in \N}$ is a Cauchy sequence in $W$. Since
  $W$ is a Banach space, the limit $\lim_{n\to \infty}A_nf$ exists and
  is an element of $W$. Thus, we can define the operator
  \begin{align*}
  A\cl V & \to W\\
  f & \mapsto \lim_{n\to \infty}A_nf,
  \end{align*}
  called the pointwise limit of $\{A_n\}_{n\in \N}$.
\item We now show that $A\in \mathcal{L}(V,W)$. This is where
  the previous lemma comes in handy. For linearity, let $f,g\in V$ and
  $z\in \C$. Then
  \begin{align*}
  A(zf+g) &:= \lim_{n\to \infty} A_n(zf+g)\\
  & = \lim_{n\to \infty} (zA_nf+A_ng)\\
  & = z\lim_{n\to \infty} A_nf+\lim_{n\to \infty}A_ng\\
  &=: zAf+Ag
  \end{align*}
  where we have used the linearity of each $A_n$ and part (i) and (ii)
  of \Cref{lem:forcompl}.
  Now we prove that \(A\) is bounded.
  First recall that $\{A_n\}_{n\in \N}$ is a Cauchy sequence. It
  follows that \(\{\|A_n\|\}_{n\in\N}\) is also Cauchy. Indeed, for
  \(\epsilon>0\) there exists $N\in \N$ such that
  \begin{equation*}
  \forall \, n,m \geq N : \ \|A_n-A_m\|<\varepsilon.
  \end{equation*}
  It follows that
  \begin{equation*}
  \bigl| \|A_n\|-\|A_m\|\bigr| \leq  \|A_n-A_m\|<\varepsilon
  \end{equation*}
  for all $n,m \geq N$.
  Since \(\R\) is complete, \(\{\|A_n\|\}_{n\in\N}\) converges to a
  finite \(r\).
  Now, by part (ii) and (iii) of \Cref{lem:forcompl}, we get
  \begin{align*}
  \|Af\|_W  & = \lim_{n\to \infty} \|A_nf\|_W\\
  & \leq \lim_{n\to \infty} \|A_n\|_W\|f\|_V \\
  & = \|f\|_V \lim_{n\to \infty} \|A_n\|
  \end{align*}
  for any $f\in V$.
  It follows that \(A\) is bounded and
  \begin{equation*}
  \|A\|\leq \lim_{n\to \infty} \|A_n\|
  \end{equation*}

\item
  To conclude, we show that $A_n\to A$.
  Let $\varepsilon >0$. 
  Since $\{A_n\}_{n\in \N}$ is Cauchy, there exists $N\in \N$ such that
  \begin{equation*}
  \forall \, n,m\geq N : \ \|A_n-A_m\| < \epsilon.
  \end{equation*}
  Hence, for any \(f\in V\), we have
  \begin{equation}
    \|A_nf-A_mf\|\leq \|A_n-A_m\|\|f\| <\epsilon\|f\|
  .\end{equation}
  Taking \(m\to\infty\), we get
  \begin{equation}
    \|A_nf-Af\|\leq \epsilon\|f\|
  \end{equation}
  for all \(f\in V\).
  This shows that \(\|A_n-A\|\leq\epsilon\). Since \(\epsilon>0\) is
  arbitrary, we conclude \(\|A_n-A\|\to 0\), so \(A_n\to A\).
\end{enumerate}
This concludes the proof that $(\mathcal{L}(V,W),\|\cdot\|)$ is a Banach space.
\end{proof}

The following is an extremely important special case of $\mathcal{L}(V,W)$.

\begin{definition}
Let $V$ be a normed space. Then $V^*:=\mathcal{L}(V,\C)$ is called the \emph{dual}\index{dual} of $V$.
\end{definition}

Note that, since $\C$ is a Banach space, the dual of a normed space is a Banach space. The elements of $V^*$ are variously called \emph{covectors} or \emph{functionals} on $V$. 
\begin{remark}
You may recall from undergraduate linear algebra that the dual of a vector space was defined to be the vector space of \emph{all} linear maps $V\to\C$, rather than just the bounded ones. This is because, in finite dimensions, all linear maps are bounded. So the two definitions agree as long as we are in finite dimensions. If we used the same definition for the infinite-dimensional case, then $V^*$ would lack some very desirable properties, such as that of being a Banach space.
\end{remark}
The dual space can be used to define a weaker notion of convergence called, rather unimaginatively, weak convergence. 

\begin{definition}
A sequence $\{f_n\}_{n\in \N}$ is said to \emph{converge
weakly}\index{weak convergence} to $f\in V$ if the sequence of complex
numbers $\{\varphi(f_n)\}_{n\in \N}$ converges to \(\varphi(f)\) for
every linear functional \(\varphi\):
\begin{equation*}
\forall \, \varphi\in V^* : \lim_{n\to\infty}\varphi(f_n) = \varphi(f).
\end{equation*}
To indicate that the sequence $\{f_n\}_{n\in \N}$ converges weakly to
$f\in V$ we write
\begin{equation*}
\wlim_{n\to \infty} f_n = f.
\end{equation*}
\end{definition}

In order to further emphasise the distinction with weak convergence,
we may say that $\{f_n\}_{n\in \N}$ converges \emph{strongly} to $f\in
V$ if it converges according to the usual definition, and we will
write accordingly
\begin{equation*}
\slim_{n\to \infty} f_n = f.
\end{equation*}

\begin{proposition}
Let $\{f_n\}_{n\in \N}$ be a sequence in a normed space $(V,\|\cdot\|_V)$. If $\{f_n\}_{n\in \N}$ converges strongly to $f\in V$, then it also converges weakly to $f\in V$, i.e.\
\begin{equation*}
\slim_{n\to \infty} f_n = f \quad \Rightarrow \quad \wlim_{n\to \infty} f_n = f.
\end{equation*}
\end{proposition}

\begin{proof}
Let $\varepsilon >0$ and let $\varphi\in V^*$. Since $\{f_n\}_{n\in \N}$ converges strongly to $f\in V$, there exists $N\in\N$ such that
\begin{equation*}
\forall \, n\geq N : \ \|f_n-f\|_V<\frac{\varepsilon}{\|\varphi\|}.
\end{equation*}
Then, since $\varphi\in V^*$ is bounded, we have
\begin{align*}
|\varphi (f_n)-\varphi (f)| & = |\varphi ( f_n-f)|\\
& \leq \|\varphi\|\|f_n-f\|_V\\
&< \|\varphi\|\frac{\varepsilon}{\|\varphi\|}\\
& = \varepsilon
\end{align*}
for any $n\geq N$. Hence, $\displaystyle  \lim_{n\to \infty} \varphi(f_n) = \varphi(f)$. That is, $\displaystyle  \wlim_{n\to \infty} f_n = f$.
\end{proof}

\section{Extension of bounded linear operators}

Note that, so far, we have only considered bounded linear maps
$A\cl\mathcal{D}_A\to W$ where $\mathcal{D}_A$ is the whole of $V$,
rather than a subspace thereof. The reason for this is that we will
only consider densely defined linear maps in general, and any bounded
linear map from a dense subspace of $V$ can be extended to a bounded
linear map from the whole of $V$. Moreover, the extension is unique.
This is the content of the so-called BLT\footnote{Bounded Linear
Transformation, \emph{not} Bacon, Lettuce, Tomato.
\begin{tikzpicture}[baseline={($ (current bounding box.center)-
(0,-8pt) $)}]\includegraphics[scale=0.12]{blt}\end{tikzpicture}}
theorem.

\begin{lemma}
Let $(V,\|\cdot\|)$ be a normed space and let $\mathcal{D}_A$ be a
dense subspace of $V$. Then, for any $f\in V$, there exists a sequence
$\{f_n\}_{n\in \N}$ in $\mathcal{D}_A$ which converges to $f$.
\end{lemma}
\begin{proof}
Let $f\in V$. Since \(\Cal D_A\) is dense in \(V\), for every
\(n\in\N\), there is an \(f_n\in\Cal D_A\) with \(\|f_n-f\|<
\frac{1}{n}\).
We obtain a sequence $\{f_n\}_{n\in \N}$ in $\mathcal{D}_A$
and, for each \(\epsilon > 0\) there is \(n\) with
\(\frac{1}{n}<\epsilon\), so
\begin{equation*}
\|f_n-f\| < \frac{1}{n} < \epsilon,
\end{equation*}
which shows that \(f_n\to f\).
\end{proof}

\begin{definition}
Let $V,W$ be vector spaces and let $A\cl\mathcal{D}_A\to W$ be a
linear map, where $\mathcal{D}_A\subseteq V$. An
\emph{extension}\index{extension} of $A$ is a linear map $\widehat
A\cl V\to W$ such that
\begin{equation*}
\forall \, \alpha\in \mathcal{D}_A : \ \widehat A \alpha = A\alpha.
\end{equation*}
\end{definition}

\begin{theorem}[BLT theorem]
  Let $V$ be a normed space and $W$ a Banach space. Any densely defined
  bounded linear map $A\cl\mathcal{D}_A\to W$ has a unique bounded extension
  $\widehat A\cl V\to W$. Moreover, $\|{\widehat A}\|=\|A\|$.
\end{theorem}

\begin{proof}
  Let $A\in \mathcal{L}(\mathcal{D}_A,W)$. Since $\mathcal{D}_A$
  is dense in $V$, for any $f\in V$ there exists a sequence
  $\{f_n\}_{n\in \N}$ in $\mathcal{D}_A$ which converges to $f$.
  Moreover, since $A$ is bounded, we have 
  \begin{equation*}
  \forall \, m,n\in \N : \ \|Af_n-Af_m\|_W \leq \|A\| \|f_n-f_m\|_V,
  \end{equation*}
  from which it quickly follows that $\{Af_n\}_{n\in \N}$ is Cauchy in
  $W$. As $W$ is a Banach space, this sequence converges to an element
  of $W$.
  Moreover, if
  $\{\alpha_n\}_{n\in \N}$ and $\{\beta_n\}_{n\in \N}$ be two
  sequences in $\mathcal{D}_A$ which converge to $f\in V$,
  Then for all \(n\) we have
  \begin{align*}
    \|A\alpha_n-A\beta_n\|_W
    & = \|A(\alpha_n-\beta_n)\|_W\\
    &\leq \|A\| \| \alpha_n-\beta_n\|_V\\
    & = \|A\| \| \alpha_n-f+f-\beta_n\|_V\\
    &\leq \|A\| (\| \alpha_n-f\|_V+\|f-\beta_n\|_V) \to 0
  \end{align*}
  where we have used the fact that $A$ is bounded.
  Thus, we have shown that \((A\alpha_n-A\beta_n)\to 0\), so
  \begin{equation*}
  \lim_{n\to\infty}A\alpha_n = \lim_{n\to\infty}A\beta_n.
  \end{equation*}
  Hence, we can define a function
  \begin{align*}
  \widehat A \cl V &\to W\\
  f & \mapsto \lim_{n\to\infty}Af_n.
  \end{align*}
  where \(\{f_n\}\) is any sequence in \(V\) converging to \(f\).
  By \Cref{lem:forcompl}, it follows easily that \(\widehat A\) is
  linear. Taking constant sequences, we see that it is an extension of
  \(A\).
  It is also bounded, since for every \(f\in V\) and every sequence
  \(f_n\to f\), we have
  \begin{align*}
  \|\widehat A f\|_W
    &:= \bigl\|\lim_{n\to\infty}Af_n\bigr\|_W\\
    &= \lim_{n\to\infty}\|Af_n\|_W\\
    & \leq\lim_{n\to\infty}\|A\|\|f_n\|_V\\
    & = \|A\|\lim_{n\to\infty}\|f_n\|_V\\
    & = \|A\|\|f\|_V,
  \end{align*}
  which also shows that \(\|\widehat A\|\leq\|A\|\).
  The other inequality \(\|A\|\leq\|\widehat A\|\) comes from the fact
  that \(\widehat A\) is an extension of \(A\).
  Finally, if \(A'\) is another bounded linear extension of \(A\), we
  have
  \begin{equation}
    A'f
    = A'(\lim_{n\to\infty}f_n)
    = \lim_{n\to\infty}A'f_n
    = \lim_{n\to\infty}Af_n
    = \lim_{n\to\infty}Af_n
    = \widehat Af
  ,\end{equation}
  since bounded maps are continuous.
\end{proof}

\begin{remark}
  The norms on each side of the equality $\|\widehat A\|=\|A\|$
  are norms over different spaces.
  A more explicit way to note this would be to write
  \begin{equation*}
  \|\widehat A\|_{\mathcal{L}(V,W)} = \|A\|_{\mathcal{L}(\mathcal{D}_A,W)},
  \end{equation*}
\end{remark}





















